%! Author = lili
%! Date = 8/2/26

\section{Alignment categories and types}
% TODO List of Alignment types
% gap-free local alignments
% free end-gap alignment
% non-overlapping suboptimal local alignment
% (sub)optimal local alignment
% local alignment
% (sub)optimal alignment




\paragraph{Global vs local}
Bei global und lokal geht es um die Art von Algorithmus, die genutzt wird bzw. welche Bereiche von Sequenzen
aligniert werden sollen. \cite{tamang}

Weitere Beschreichsabgrenzungen (oder so...) wären noch:
\begin{itemize}
    \item Semi-global alignment (/ free-end-gap alignment)
\end{itemize}

\paragraph{Pairwise vs Multiple Sequence}
Hierbei geht es um den Alignment-TYP, bzw. wie viele Sequenzen aligned werden sollen. \cite{tamang}

\paragraph{Optimal vs. Suboptimal}
Hierbei geht es um die Qualität des alignments (in bezug auf das Scoring-Modell!)

\paragraph{free end-gap und so} Dabwei geht es um die Gap-Kosten bzw. wie mit welchen Arten von Gaps umgegangen
werden soll.
\begin{itemize}
    \item Gap-free alignment
    \item Linear gap penalty
    \item Affine gap penalty
    \item Free end gaps / end-gap-free
\end{itemize}

\paragraph{Overlapping vs non-overlapping}
Das ist letztendlich eine weitere Beschränkung unter dem größeren Schirm davon Alternativen/Mehrere Alignments zu
finden.

In ähnlicher Motivation könnte man nooch listen:
\begin{itemize}
    \item k-best alignments
    \item multiple optimal alignments
\end{itemize}



